% Template: Góc AOB với hai tia
% Dùng cho: đo góc, tia phân giác, góc kề bù, góc đối đỉnh
% --------------------------------------------------------
\documentclass[border=5pt]{standalone}
\usepackage[utf8]{inputenc}
\usepackage[T5]{fontenc}
\usepackage[vietnamese]{babel}
\usepackage{tikz}
\usepackage{helvet}
\renewcommand{\familydefault}{\sfdefault}
% ── Hệ màu tập trung (chỉnh tại đây để đổi theme) ──────────
\usepackage{xcolor}
\definecolor{mauchinh} {HTML}{1565C0}
\definecolor{mauphu}   {HTML}{43A047}
\definecolor{maunoibat}{HTML}{E53935}
\definecolor{mauvien}  {HTML}{424242}
{mauchinh} {HTML}{1565C0}
{mauphu}   {HTML}{43A047}
{maunoibat}{HTML}{E53935}
{mauvien}  {HTML}{424242}
\usetikzlibrary{calc, angles, quotes, arrows.meta}

\begin{document}
\begin{tikzpicture}[
    scale=1.2,
    >=Stealth,
    line join=round,
    line cap=round,
    every node/.style={font=\sffamily\small},
    thick,
]
    % Độ dài tia
    \pgfmathsetmacro{\raylen}{2.5}
    % Góc AOB = 55°
    \pgfmathsetmacro{\angAOB}{55}

    % Gốc góc O
    \path (0, 0) coordinate (O);
    % Tia OA nằm ngang
    \path (\raylen, 0) coordinate (OA);
    % Tia OB tại góc 55° so với OA
    \path (\angAOB:\raylen) coordinate (OB);

    % Tia OA
    \draw[->] (O) -- (OA);
    % Tia OB
    \draw[->] (O) -- (OB);

    % Cung góc alpha với nhãn
    \pic[draw, angle radius=0.7cm,
         angle eccentricity=1.4,
         "$\alpha$"] {angle = OA--O--OB};

    % Nhãn gốc O
    \path (O)  node[below left] {$O$};
    % Nhãn điểm A
    \path (OA) node[right] {$A$};
    % Nhãn điểm B
    \path (OB) node[above] {$B$};

    % Điểm tại gốc O
    \fill (O) circle (2pt);
\end{tikzpicture}
\end{document}
